\documentclass[a4paper, 12pt]{article}

\usepackage[portuges]{babel}
\usepackage[utf8]{inputenc}
\usepackage{amsmath}
\usepackage{indentfirst}
\usepackage{graphicx}
\usepackage{multicol,lipsum}
\usepackage[bottom=2cm,top=2cm,left=2cm,right=2cm]{geometry}

\begin{document}
%\maketitle

\begin{titlepage}
	\begin{center}
	
	%\begin{figure}[!ht]
	%\centering
	 % \includegraphics[width=4cm]{C:/ifba.jpg}
	%\end{figure}

		\Large{Instituto Federal de Educação Tecnológica da Bahia}\\
		\large{Departamento de Ensino}\\ 
		\large{Coordenação de Engenharia Elétrica}\\ 
        \large{Disciplina Medidas Elétricas e Magnéticas}\\ 
        \vspace{115pt}
        \textbf{\Large{Aluno(a) fulano de tal da Silva }}\\
		\vspace{15pt}
        \vspace{95pt}
        \textbf{\Large{Experimento 1: (título) }}\\
		%\title{{\large{Título}}}
		\vspace{3,5cm}
	\end{center}
	
	%\begin{flushleft}
	%	\begin{tabbing}
	%		Professor orientador: \\
	%		Professor co-orientador: \\
	%\end{tabbing}
 %\end{flushleft}
	\vspace{1cm}
	
	\begin{center}
		\vspace{\fill}
			 mês\\
		 ano
			\end{center}
\end{titlepage}
%%%%%%%%%%%%%%%%%%%%%%%%%%%%%%%%%%%%%%%%%%%%%%%%%%%%%%%%%%%

% % % % % % % % %FOLHA DE ROSTO % % % % % % % % % %

\begin{titlepage}
	\begin{center}
	
	%\begin{figure}[!ht]
	%\centering
	 % \includegraphics[width=4cm]{C:/ifba.jpg}
	%\end{figure}

		\Large{Instituto Federal de Educação Tecnológica da Bahia}\\
		\large{Departamento de Ensino}\\ 
		\large{Coordenação de Engenharia Elétrica}\\ 
        \large{Disciplina Medidas Elétricas e Magnéticas}\\ 
        \vspace{115pt}
        \textbf{\Large{Nome do Aluno }}\\
		\vspace{15pt}
        \vspace{95pt}
        \textbf{\Large{Experimento 1: (título) }}\\
		%\title{{\large{Título}}}
		\vspace{3,5cm}
	\end{center}
	
	\begin{flushright}

   \begin{list}{}{
      \setlength{\leftmargin}{10cm}
      \setlength{\rightmargin}{0cm}
      \setlength{\labelwidth}{0pt}
      \setlength{\labelsep}{\leftmargin}}

      \item Relatório da disciplina Medidas Elétricas e Magnéticas do Curso de Engenharia Elética do IFBA Campus Paulo Afonso, como requisito parcial para avaliação .

      \begin{list}{}{
      \setlength{\leftmargin}{0cm}
      \setlength{\rightmargin}{0cm}
      \setlength{\labelwidth}{0pt}
      \setlength{\labelsep}{\leftmargin}}

			
            \item Aluno(a): \
      		\item Professor da disciplina: \

      \end{list}
   \end{list}
\end{flushright}
\vspace{1cm}
\begin{center}
		\vspace{\fill}
		 mês\\
		 ano
			\end{center}
\end{titlepage}

ewpage
% % % % % % % % % % % % % % % % % % % % % % % % % % % % % % % % % % % % %

ewpage
\tableofcontents
\thispagestyle{empty}


ewpage
\pagenumbering{arabic}
% % % % % % % % % % % % % % % % % % % % % % % % % % %
\section{Introdução}

ewpage
\section{Objetivos}

ewpage

\section{Material utilizado}


ewpage
\section{Procedimentos Experimentais}

ewpage
\section{Resultados e discussões}

ewpage

\section{Conclusões}

ewpage

\addcontentsline{toc}{section}{Referências Bibliográficas}
\section*{Referências Bibliográficas}
\footnotesize{


oindent AGUIRRE, L. A. Introdução à Identificação de Sistemas, Técnicas Lineares e Não lineares Aplicadas a Sistemas Reais. Belo Horizonte, Brasil, EDUFMG. 2004.(exemplo)\\

}

ewpage
\addcontentsline{toc}{section}{Anexo}
\section*{Anexo}
\end{document}



